%% paper/sections/related_work.tex
%% Draft only. Literature search for prior leakage reports (OUC-CGE specifically,
%% and clip-derived affect/engagement benchmarks generally) was run BEFORE this
%% file was drafted, per instruction. Summary: no prior report of OUC-CGE leakage
%% was found; a directly analogous, previously-reported case (ChaLearn First
%% Impressions) and a contemporaneous general leakage-detection methodology were
%% found and are cited below, so \S\ref{sec:dataset-audit} is framed as
%% independent confirmation of a known failure class plus a reusable test, not a
%% novelty claim.
\label{sec:related-work}

\paragraph{Efficient video attention.}
Reducing the cost of full space-time self-attention over a patch grid has
been approached by factorising it (divided space-time
attention~\cite{bertasius2021timesformer}, tubelet and factorised encoders
in ViViT~\cite{arnab2021vivit}), by restricting its span (windowed and
pooling attention in Multiscale Vision
Transformers~\cite{fan2021multiscale}), by merging redundant tokens
post hoc~\cite{bolya2023tome}, and by replacing attention's quadratic cost
with a state-space recurrence~\cite{gu2023mamba}. These methods are
orthogonal to the question this paper asks: they make attention over the
same referent-free grid cheaper or shorter-range, but the tokens being
attended over still lack a persistent identity across time. EPT does not
propose a cheaper approximation of grid attention --- it proposes a
different definition of what a token denotes, orthogonal to (and in
principle combinable with) any of the above.

\paragraph{Object- and entity-centric video representation.}
Slot Attention~\cite{locatello2020slot} and its video extension
SAVi~\cite{kipf2022savi} learn object-centric representations end-to-end,
discovering slots without external detection or tracking supervision; they
motivate the persistence hypothesis this paper tests but do not isolate
persistence from token count or spatial localization as separate causal
factors, and typically target scene decomposition or dynamics prediction
rather than a downstream classification task with a matched-budget grid
baseline. EPT takes the more constrained, less general path deliberately:
identity comes from a detector and tracker rather than being discovered, in
exchange for a token stream where identity persistence is a manipulable,
independently-controllable variable (\S\ref{sec:method}) rather than an
emergent property that is hard to ablate.

\paragraph{Engagement and group affect recognition.}
Video-based engagement recognition has largely developed on datasets built
around a single subject and a subject-disjoint split --- DAiSEE, 9068 ten
second e-learning clips from 112 subjects, four ordinal engagement
levels~\cite{gupta2016daisee}; EngageNet's official split is likewise
subject-independent. Multi-party affect work (e.g.\ MELD, dialogue-level
sentiment and emotion over multi-speaker television
video~\cite{poria2019meld}) instead splits by dialogue or scene. OUC-CGE
extends the task to \emph{group} engagement in real classrooms, released
under a split described as 80/10/10 by randomly generated clip
filename~\cite{lu2025ouccge}. Neither the release nor, to our search, any
subsequent use of it discusses whether that random clip-level assignment
keeps physical recording sessions disjoint across splits;
\S\ref{sec:dataset-audit} finds that it does not.

\paragraph{Split leakage in clip-derived video benchmarks.}
This is not the first video benchmark found to leak across splits at the
clip level. In the ChaLearn First Impressions apparent-personality corpus
--- 15-second clips cut from a smaller pool of longer YouTube videos ---
overlapping data splits are identified as a confound in the
personality-attribution literature, with the same source video contributing
clips to more than one split~\cite{ras2019background}. The
general problem of detecting such leakage in visual datasets via embedding
similarity and cluster/connected-component analysis is addressed directly,
contemporaneously with this project's own audit, by
Ramos~et~al.~\cite{ramos2025leakage}, using a similar toolkit (embedding
nearest-neighbor similarity, connected components across splits) to the one
independently developed in \S\ref{sec:dataset-audit}. We did not find a
prior report specific to OUC-CGE --- it is a recent release
(April 2025~\cite{lu2025ouccge}) with, to our search, no published critique
or errata. We therefore present \S\ref{sec:dataset-audit}'s finding as an
instance of an already-documented failure class in clip-derived video
benchmarks, newly identified for this specific dataset, rather than as a
novel phenomenon; the contribution is the finding for OUC-CGE plus a
reusable, dataset-agnostic admission procedure, not the general observation
that clip-level splitting can leak.

\paragraph{Edge video analytics.}
Deploying video models under CPU/edge power and latency budgets typically
trades accuracy for cheaper backbones or lower frame rates. This paper
reports cost as an end-to-end pipeline number, detection included, rather
than as attention FLOPs alone, following the general caution in edge
deployment work that a component-level efficiency gain does not imply a
system-level one (\S\ref{sec:results}).
